\documentclass[12pt,letterpaper]{article}
\usepackage[margin=2.4cm, letterpaper]{geometry}

\usepackage[utf8]{inputenc}
\usepackage[spanish]{babel}
\usepackage{times} 
\usepackage{setspace}
\usepackage{hanging}

\setlength{\parindent}{0pt}       
\onehalfspacing                  
\usepackage{titlesec}


\titleformat{\section}
  {\normalfont\large\bfseries\raggedright}
  {\thesection}{1em}{}
\titlespacing*{\section}{0pt}{3.5ex plus 1ex minus .2ex}{1ex}

\usepackage{hyperref}
\usepackage{cite}
\usepackage{xurl}
\begin{document}

\begin{center}
    {\large\textbf{INSTRUMENTO DE RECOLECCION DE INFORMACION}} \\[1.5ex]
    \textit{J. J. CHÁN CUELLAR } \\[1.5ex]
    7690-22-1805 \quad Universidad Mariano Gálvez \\[1.0ex]
    Seminario de tecnologías de información \\[1.0ex]
    \textit{jchanc4@miumg.edu.gt}
\end{center}
\vspace{1em}
\section*{Resumen}
El presente ensayo busca dar a conocer aspectos importantes sobre cómo realizar una encuesta efectiva y óptima, identificando los elementos necesarios para diseñarla correctamente y obtener información confiable, clara y útil para la toma de decisiones. Una encuesta correctamente elaborada permite recopilar datos relevantes de una población determinada, siempre que el propósito de la investigación esté bien establecido y que las preguntas sean seleccionadas y redactadas de manera adecuada para evitar errores que puedan afectar los resultados. Para la elaboración de este trabajo se investigaron diferentes aspectos relacionados con la metodología de encuestas, el diseño de cuestionarios, selección de muestras, tipos de preguntas y buenas prácticas para mejorar la calidad de la información obtenida. Como resultado del análisis realizado, se determinó que la efectividad de una encuesta depende de una planificación adecuada, donde es necesario definir claramente lo que se desea medir, aplicar preguntas comprensibles y realizar una revisión previa del instrumento antes de su aplicación. Se concluye que una encuesta no consiste únicamente en formular preguntas, sino en aplicar un proceso ordenado que permita obtener datos válidos y representativos, reduciendo posibles errores y facilitando un análisis más confiable de la información recolectada.


\vspace{1em}

\noindent\textbf{Palabras claves:} encuesta, cuestionario, metodología, diseño de preguntas, recolección de información
\vspace{1em}

\section*{Desarrollo del tema}
Una encuesta es una técnica de investigación utilizada para obtener información de una población determinada mediante la aplicación de un conjunto de preguntas estructuradas. Su finalidad principal es recopilar datos que permitan conocer opiniones, características, comportamientos o situaciones específicas relacionadas con un fenomeno de estudio. Para que una encuesta sea considerada efectiva y óptima no basta únicamente con formular preguntas y obtener respuestas, sino que requiere de una planificación adecuada donde cada elemento del instrumento tenga una relación directa con la información que se desea obtener.

El primer aspecto fundamental para realizar una encuesta correctamente es establecer con claridad qué información se desea recopilar y cuál será la utilidad de los resultados obtenidos. Antes de iniciar la elaboración del cuestionario es necesario definir el problema de investigación y los objetivos que orientan el proceso, debido a que estos elementos permiten determinar qué preguntas son necesarias y cuáles pueden generar información innecesaria. Una encuesta que no tiene una finalidad claramente establecida puede producir datos difíciles de interpretar y resultados que no aporten soluciones al problema planteado.

Es importante corresponder a la selección de la población y la muestra. La población representa el conjunto total de personas o elementos que poseen características relacionadas con la investigación, mientras que la muestra corresponde a una parte seleccionada de esa población que será analizada. La elección adecuada de la muestra permite obtener resultados más representativos y disminuir errores en la interpretación de la información. En investigaciones donde se busca generalizar los resultados, es recomendable utilizar métodos de selección probabilísticos; En determinados estudios también pueden utilizarse muestras no probabilísticas, siempre considerando sus limitaciones y el alcance de las conclusiones obtenidas.

El diseño del cuestionario representa uno de los puntos más importantes dentro del desarrollo de una encuesta, debido a que la calidad de las preguntas influye directamente en la calidad de los datos recolectados. Las preguntas deben estar redactadas utilizando un lenguaje claro, sencillo y comprensible para los participantes, evitando términos técnicos que puedan generar confusión. Además, deben estar enfocadas en obtener información específica y evitar preguntas que incluyan opiniones del investigador o que puedan influenciar la respuesta del encuestado.

Dentro de la elaboración de preguntas existen diferentes tipos de estructuras que pueden utilizarse dependiendo de la información que se desea obtener. Las preguntas cerradas permiten seleccionar respuestas previamente establecidas y facilitan el análisis estadístico de los resultados, mientras que las preguntas abiertas permiten que los participantes expresen sus opiniones con mayor libertad, proporcionando información más detallada. La elección entre ambos tipos dependerá de los objetivos del estudio y del tipo de análisis que posteriormente se realizará.

Un aspecto que debe considerarse durante el diseño de una encuesta es la reducción del error de respuesta. Este problema ocurre cuando las respuestas obtenidas no representan completamente la realidad debido a factores relacionados con la forma en que se realizó la pregunta, la interpretación del participante o incluso la influencia del contexto. Para disminuir este riesgo es recomendable utilizar preguntas neutrales, evitar expresiones que orienten una respuesta específica y proporcionar opciones equilibradas que permitan al participante responder según su verdadera opinión.

Es importante considerar la organización y estructura del cuestionario. El orden de las preguntas puede influir en la disposición del participante para completar la encuesta, por lo que generalmente se recomienda iniciar con preguntas sencillas que permitan familiarizar al encuestado con el instrumento. Las preguntas más complejas o sensibles pueden colocarse posteriormente, mientras que la información personal o datos de clasificación suelen ubicarse al final para evitar generar rechazo desde el inicio.

Antes de aplicar una encuesta a toda la población seleccionada es recomendable realizar una prueba piloto. Esta etapa consiste en aplicar el instrumento a un grupo reducido de personas con características similares a la población objetivo para identificar posibles errores en la redacción, problemas de comprensión o dificultades técnicas. La prueba piloto permite realizar mejoras antes de la aplicación definitiva y aumenta la posibilidad de obtener información más precisa.

En la actualidad, las herramientas digitales han facilitado considerablemente el proceso de elaboración y aplicación de encuestas. Plataformas como Google Forms permiten crear cuestionarios en línea, recopilar respuestas automáticamente y organizar los datos para su posterior análisis. Sin embargo, aunque las herramientas tecnológicas facilitan el proceso, no sustituyen la importancia de una correcta planificación metodológica, ya que una encuesta mal diseñada continuará generando información poco confiable aunque sea aplicada mediante una plataforma avanzada.

El análisis de los datos obtenidos representa una etapa fundamental para determinar el valor de la encuesta realizada. La información recolectada debe ser organizada, revisada e interpretada considerando los objetivos establecidos inicialmente. Además, es importante aplicar principios éticos durante todo el proceso, protegiendo la privacidad de los participantes y utilizando la información únicamente para los fines establecidos en la investigación.
\vspace{1em}

\section*{Observaciones y comentarios}
Durante la investigación realizada logre identificar que la efectividad de una encuesta depende principalmente de la calidad del proceso metodológico utilizado para su elaboración y no únicamente de la herramienta empleada para aplicarla.También se determinó que la reducción de errores y la correcta formulación de preguntas representan factores fundamentales para obtener resultados más precisos y útiles en una investigación.
\vspace{1em}

\section*{Conclusiones}
\begin{enumerate}
    \item La elaboración de una encuesta efectiva requiere una planificación adecuada, donde se establezca claramente la información que se desea obtener y la forma correcta de recopilarla.
    \item El uso de herramientas digitales facilita la aplicación y análisis de encuestas, pero no sustituye la importancia de aplicar una metodología correcta durante todo el proceso de investigación.
    \item una encuesta mal formulada no genera información útil ni confiable, por lo que estoy de acuerdo en que la calidad de las preguntas influye directamente en la calidad de los resultados obtenidos.
    \item El punto de una encuesta es dar informacion pero depende de la calidad de informacion nos puede ayudar o nos puede afectar a nuestra investigacion.
\end{enumerate}
\vspace{1em}
\section*{7. Bibliografía}

\hangindent=1.27cm

Meneses, J. (2016). 
\textit{El cuestionario}. 
Universitat Oberta de Catalunya. 
https://femrecerca.cat/meneses/publication/cuestionario/cuestionario.pdf

\hangindent=1.27cm

Universidad de Granada. (s. f.). 
\textit{Tema 2: El cuestionario. Diseño del cuestionario}. 
Departamento de Estadística e Investigación Operativa. 
https://www.ugr.es/~diploeio/documentos/tema2.pdf

\hangindent=1.27cm

Universidad Veracruzana. (2024). 
\textit{El cuestionario como instrumento de investigación/evaluación}. 
https://lumen.uv.mx/resources/files/documents/2024/2/1/9868/4c5004ac-944c-4e29-a170-b6d4b10907cd.pdf

\hangindent=1.27cm

Instituto de Derechos Humanos de la Universidad Nacional de La Plata. (s. f.). 
\textit{Metodología de la investigación}. 
Universidad Nacional de La Plata. 
http://www.derechoshumanos.unlp.edu.ar/assets/files/documentos/metodologia-de-la-investigacion.pdf

\section*{Encuesta elaborada}

El instrumento de recolección de información fue elaborado mediante Google Forms.
La encuesta puede ser consultada en el siguiente enlace:

\href{https://docs.google.com/forms/d/e/1FAIpQLSdn-KN6MWQF3YVcPPS4IwZZOUlooedp1UIy0hJ0T31uJIjDgg/viewform?usp=publish-editor}
{Encuesta sobre el conocimiento y uso de la Inteligencia Artificial en estudiantes}.
\end{document}